\documentclass[11pt]{article}

% Shared notes settings for Elements of AIML lectures (UPES).
% Keep this file free of \documentclass to allow reuse via \input.

\usepackage[utf8]{inputenc}
\usepackage[T1]{fontenc}
\usepackage[a4paper,margin=1in]{geometry}
\usepackage{hyperref}
\usepackage{xurl}
\usepackage{booktabs}
\usepackage{amsmath}
\usepackage{amssymb}
\usepackage{listings}
\usepackage{xcolor}
\usepackage{enumitem}
\usepackage{parskip}

% Code listings (general-purpose).
\lstset{
  basicstyle=\ttfamily\small,
  keywordstyle=\color{blue},
  commentstyle=\color{gray},
  stringstyle=\color{teal},
  showstringspaces=false,
  columns=fullflexible,
  frame=single,
  framerule=0.2pt,
  breaklines=true
}

\setlist[itemize]{topsep=0.3em,itemsep=0.2em}
\setlist[enumerate]{topsep=0.3em,itemsep=0.2em}

% Simple per-section source line.
\newcommand{\notesources}[1]{\par\noindent{\small\color{gray}\textit{Sources:} \sloppy #1\par}}

% Unit-wise source strings for this subject (used by slides/notes).
% Path is relative to the lecture `latex/` folder (the compilation CWD).
\IfFileExists{../../../_shared/aiml-sources.tex}{%
  % Source strings for Elements of AIML (used by slides + notes).

\newcommand{\AIMLRepoURL}{https://github.com/tali7c/Elements-of-AIML}

\newcommand{\AIMLLabSources}{%
Provost and Fawcett, \textit{Data Science for Business}.;%
McKinney, \textit{Python for Data Analysis}.;%
WEKA Documentation (Waikato Environment for Knowledge Analysis), accessed 2026-02-28.;%
Matplotlib and Seaborn documentation, accessed 2026-02-28.%
}

\newcommand{\AIMLUnitISources}{%
Rich and Knight, \textit{Artificial Intelligence}, McGraw Hill, 2017.;%
Russell and Norvig, \textit{Artificial Intelligence: A Modern Approach} (4e), Ch. 1--2 (Introduction, Intelligent Agents).;%
Poole and Mackworth, \textit{Artificial Intelligence: Foundations of Computational Agents} (2e), selected sections.%
}

\newcommand{\AIMLUnitIISources}{%
Russell and Norvig, \textit{Artificial Intelligence: A Modern Approach} (4e), Ch. 7--9 (Logical Agents, First-Order Logic, Inference).;%
Huth and Ryan, \textit{Logic in Computer Science} (2e), selected sections on propositional and first-order logic.;%
Genesereth and Nilsson, \textit{Logical Foundations of Artificial Intelligence}, selected chapters.%
}

\newcommand{\AIMLUnitIIISources}{%
Mueller and Massaron, \textit{Machine Learning for Dummies}, For Dummies, 2016.;%
Mitchell, \textit{Machine Learning}, selected chapters on learning basics and evaluation.;%
Scikit-learn documentation, accessed 2026-02-28.%
}

\newcommand{\AIMLUnitIVSources}{%
Mueller and Massaron, \textit{Machine Learning for Dummies}, For Dummies, 2016.;%
James, Witten, Hastie, and Tibshirani, \textit{An Introduction to Statistical Learning} (2e), selected chapters.;%
Scikit-learn documentation, accessed 2026-02-28.%
}

\newcommand{\AIMLUnitVSources}{%
NIST, \textit{AI Risk Management Framework (AI RMF 1.0)}, 2023.;%
Barocas, Hardt, and Narayanan, \textit{Fairness and Machine Learning} (online book), selected sections.;%
Knaflic, \textit{Storytelling with Data}, selected chapters on communicating results.%
}
%
}{%
  % Faculty copies live under `private/` so their compile CWD is different.
  \IfFileExists{../../../../public/_shared/aiml-sources.tex}{%
    % Source strings for Elements of AIML (used by slides + notes).

\newcommand{\AIMLRepoURL}{https://github.com/tali7c/Elements-of-AIML}

\newcommand{\AIMLLabSources}{%
Provost and Fawcett, \textit{Data Science for Business}.;%
McKinney, \textit{Python for Data Analysis}.;%
WEKA Documentation (Waikato Environment for Knowledge Analysis), accessed 2026-02-28.;%
Matplotlib and Seaborn documentation, accessed 2026-02-28.%
}

\newcommand{\AIMLUnitISources}{%
Rich and Knight, \textit{Artificial Intelligence}, McGraw Hill, 2017.;%
Russell and Norvig, \textit{Artificial Intelligence: A Modern Approach} (4e), Ch. 1--2 (Introduction, Intelligent Agents).;%
Poole and Mackworth, \textit{Artificial Intelligence: Foundations of Computational Agents} (2e), selected sections.%
}

\newcommand{\AIMLUnitIISources}{%
Russell and Norvig, \textit{Artificial Intelligence: A Modern Approach} (4e), Ch. 7--9 (Logical Agents, First-Order Logic, Inference).;%
Huth and Ryan, \textit{Logic in Computer Science} (2e), selected sections on propositional and first-order logic.;%
Genesereth and Nilsson, \textit{Logical Foundations of Artificial Intelligence}, selected chapters.%
}

\newcommand{\AIMLUnitIIISources}{%
Mueller and Massaron, \textit{Machine Learning for Dummies}, For Dummies, 2016.;%
Mitchell, \textit{Machine Learning}, selected chapters on learning basics and evaluation.;%
Scikit-learn documentation, accessed 2026-02-28.%
}

\newcommand{\AIMLUnitIVSources}{%
Mueller and Massaron, \textit{Machine Learning for Dummies}, For Dummies, 2016.;%
James, Witten, Hastie, and Tibshirani, \textit{An Introduction to Statistical Learning} (2e), selected chapters.;%
Scikit-learn documentation, accessed 2026-02-28.%
}

\newcommand{\AIMLUnitVSources}{%
NIST, \textit{AI Risk Management Framework (AI RMF 1.0)}, 2023.;%
Barocas, Hardt, and Narayanan, \textit{Fairness and Machine Learning} (online book), selected sections.;%
Knaflic, \textit{Storytelling with Data}, selected chapters on communicating results.%
}
%
  }{%
    % Source strings for Elements of AIML (used by slides + notes).

\newcommand{\AIMLRepoURL}{https://github.com/tali7c/Elements-of-AIML}

\newcommand{\AIMLLabSources}{%
Provost and Fawcett, \textit{Data Science for Business}.;%
McKinney, \textit{Python for Data Analysis}.;%
WEKA Documentation (Waikato Environment for Knowledge Analysis), accessed 2026-02-28.;%
Matplotlib and Seaborn documentation, accessed 2026-02-28.%
}

\newcommand{\AIMLUnitISources}{%
Rich and Knight, \textit{Artificial Intelligence}, McGraw Hill, 2017.;%
Russell and Norvig, \textit{Artificial Intelligence: A Modern Approach} (4e), Ch. 1--2 (Introduction, Intelligent Agents).;%
Poole and Mackworth, \textit{Artificial Intelligence: Foundations of Computational Agents} (2e), selected sections.%
}

\newcommand{\AIMLUnitIISources}{%
Russell and Norvig, \textit{Artificial Intelligence: A Modern Approach} (4e), Ch. 7--9 (Logical Agents, First-Order Logic, Inference).;%
Huth and Ryan, \textit{Logic in Computer Science} (2e), selected sections on propositional and first-order logic.;%
Genesereth and Nilsson, \textit{Logical Foundations of Artificial Intelligence}, selected chapters.%
}

\newcommand{\AIMLUnitIIISources}{%
Mueller and Massaron, \textit{Machine Learning for Dummies}, For Dummies, 2016.;%
Mitchell, \textit{Machine Learning}, selected chapters on learning basics and evaluation.;%
Scikit-learn documentation, accessed 2026-02-28.%
}

\newcommand{\AIMLUnitIVSources}{%
Mueller and Massaron, \textit{Machine Learning for Dummies}, For Dummies, 2016.;%
James, Witten, Hastie, and Tibshirani, \textit{An Introduction to Statistical Learning} (2e), selected chapters.;%
Scikit-learn documentation, accessed 2026-02-28.%
}

\newcommand{\AIMLUnitVSources}{%
NIST, \textit{AI Risk Management Framework (AI RMF 1.0)}, 2023.;%
Barocas, Hardt, and Narayanan, \textit{Fairness and Machine Learning} (online book), selected sections.;%
Knaflic, \textit{Storytelling with Data}, selected chapters on communicating results.%
}
%
  }%
}


\title{Elements of AIML Lab\\Experiment 05 (After-submission Reference): EDA using Python}
\author{Tofik Ali}
\date{\today}

\begin{document}
\maketitle
\tableofcontents

\section{How to use this reference}
This reference is meant to be shared \textbf{after you submit} Experiment~05.
It demonstrates a complete EDA flow and, more importantly, shows \textbf{how to write observations} (EDA is not only plots).

\section{EDA in one sentence}
EDA answers: \emph{``What does this dataset look like, what can go wrong, and what should we do next?''}

\section{Recommended EDA sequence (teaching order)}
\subsection{1) Load + understand structure}
\begin{lstlisting}[language=Python]
import pandas as pd

df = pd.read_csv("processed.csv")   # use Exp-04 output if available
print(df.shape)
print(df.head())
print(df.info())
\end{lstlisting}
\textbf{What to write:} mention rows/columns and the main column types (numeric/categorical/date).

\subsection{2) Data quality checks (missing + duplicates)}
\begin{lstlisting}[language=Python]
miss = df.isna().sum().sort_values(ascending=False)
print(miss[miss > 0])
print("Duplicates:", df.duplicated().sum())
\end{lstlisting}
\textbf{What to write:} name the top 2--3 columns with missing values and their counts/percentages.

\subsection{3) Descriptive summaries}
\begin{lstlisting}[language=Python]
print(df.describe())
print(df.describe(include="object"))
\end{lstlisting}
\textbf{What to look for:} min/max surprises, skewed distributions, rare categories.

\subsection{4) Univariate plots (distribution of one feature)}
\begin{lstlisting}[language=Python]
import matplotlib.pyplot as plt
import seaborn as sns

sns.histplot(df["num_col"], kde=True)
plt.title("Distribution of num_col")
plt.show()

sns.boxplot(x=df["num_col"])
plt.title("Boxplot of num_col")
plt.show()
\end{lstlisting}
\textbf{Observation examples:}
\begin{itemize}[leftmargin=*]
  \item ``\texttt{num\_col} is right-skewed; most values are below 50 with a long tail up to 300.''
  \item ``The boxplot shows extreme values; we may need outlier handling or log transform.''
\end{itemize}

\subsection{5) Relationship plots (feature vs feature)}
\begin{lstlisting}[language=Python]
corr = df.select_dtypes("number").corr()
sns.heatmap(corr, cmap="coolwarm", center=0)
plt.title("Correlation heatmap (numeric features)")
plt.show()
\end{lstlisting}
\textbf{How to interpret:}
\begin{itemize}[leftmargin=*]
  \item strong positive correlation $\approx$ features increase together,
  \item strong negative correlation $\approx$ one increases as the other decreases,
  \item near zero correlation does not guarantee independence (could be non-linear).
\end{itemize}

\subsection{6) If there is a target column (optional but valuable)}
\begin{itemize}[leftmargin=*]
  \item Numeric target: scatter plot feature vs target.
  \item Class label: compare distributions by class (boxplots or histograms per class).
\end{itemize}

\section{Minimum plot set (good standard)}
\begin{itemize}[leftmargin=*]
  \item One histogram (distribution) + written comment.
  \item One boxplot (outliers/spread) + written comment.
  \item One relationship plot (scatter or heatmap) + written comment.
\end{itemize}

\section{Common EDA writing pattern (use this in reports)}
For each plot/check, write:
\begin{enumerate}[leftmargin=*]
  \item \textbf{Evidence:} what you saw (numbers/plot pattern).
  \item \textbf{Meaning:} what it suggests about the data.
  \item \textbf{Action:} what you would do next (missing handling, scaling, encoding, etc.).
\end{enumerate}

\section{Troubleshooting}
\begin{itemize}[leftmargin=*]
  \item \textbf{Plots look empty:} check column names, data types, and missing values.
  \item \textbf{Too many categories:} plot top-$k$ value counts and group rare categories as ``Other''.
  \item \textbf{Time columns:} convert with \texttt{pd.to\_datetime} before time-based plots.
\end{itemize}

\notesources{\AIMLLabSources}
\end{document}

